Write the side lengths as $b-d,b,b+d$, with $BC=b$. The semiperimeter is $s=\dfrac{3b}{2}$. If $h$ is the altitude from $A$ and $r$ the inradius, the area identities give $\dfrac{1}{2}bh=rs=\dfrac{3}{2}br$. Hence $\boxed{h:r=3:1}$.
\begin{asy}
import graph;
usepackage("amsmath");
import geometry;
size(7cm);
pair[] p={(0,4),(0,0),(5,0)}; pair labelCentre=(0,0);for(pair labelVertex:p)labelCentre+=labelVertex;labelCentre/=p.length; draw(p[0]--p[1]--p[2]--cycle); dot(p[0]);label("$A$",p[0],unit(p[0]-labelCentre));dot(p[1]);label("$B$",p[1],unit(p[1]-labelCentre));dot(p[2]);label("$C$",p[2],unit(p[2]-labelCentre));draw((0,4)--(0,0),linewidth(1.2));label("$h$",(0,2),(-1,0));
\end{asy}

